\documentclass{article}
\usepackage[utf8]{inputenc}
\usepackage{geometry}
\usepackage{listings}
\usepackage{xcolor}
\geometry{margin=1in}

\lstset{
    basicstyle=\ttfamily\small,
    breaklines=true,
    frame=single,
    backgroundcolor=\color{gray!10},
    columns=flexible
}

\lstdefinestyle{java}{
    language=Java,
    keywordstyle=\color{blue}\bfseries,
    commentstyle=\color{green!60!black},
    stringstyle=\color{red}
}

% Define document metadata in the preamble
\title{\bfseries Reviewing a Java code snippet}
\author{Souvik Sarkar \\ \texttt{sounix000@gmail.com}} % Applied your signature author block
\date{November 9, 2020}

\begin{document}

\maketitle % Generates your standard title, author, and email block

\section*{API reference}

For this section, I assumed that the method in the code sample is part of a class \texttt{NeedlesInHaystack}.

\subsection*{Class: \texttt{NeedlesInHaystack}}

This class provides functionality to search for multiple needle strings within a single haystack string and outputs the frequency of each needle's occurrence.

\subsubsection*{Method: \texttt{findNeedles}}

\begin{itemize}
    \item \textbf{Type}: \texttt{public static}
    
    \item \textbf{Parameters}:
    \begin{itemize}
        \item \texttt{String haystack}: The text within which to search for the needles.
        \item \texttt{String[] needles}: An array of strings representing the words (needles) to be counted in the haystack.
    \end{itemize}
    
    \item \textbf{Returns}: \texttt{void} (prints results to the console)
    
    \item \textbf{Usage}: This method prints the number of times each string in the \texttt{needles} array appears in the \texttt{haystack} string.
\end{itemize}

\subsubsection*{Example}

\begin{lstlisting}[style=java]
String haystack = "Google Cloud provides APIs to use Google's ML/AI capabilities.";
String[] needles = {"Google", "API", "documentation", "AWS", "ML/AI"};
new NeedlesInHaystack().findNeedles(haystack, needles);
\end{lstlisting}

\subsubsection*{Output}

\begin{lstlisting}
Google: 2
API: 0
documentation: 0
AWS: 0
ML/AI: 1
\end{lstlisting}

\section*{Suggestions for code improvement}

After reviewing the code sample, I have the following suggestions for improvement:

\begin{itemize}
    \item \textbf{Limit the number of \texttt{needles}}: The current code restricts the length of the \texttt{needles} array to five. If this is a strict requirement, modify the message in the print statement within the \texttt{if} block to ``Use a maximum of five words!''. This provides clearer guidance than ``Too many...''.
    
    \item \textbf{Assign \texttt{needles.length} to a variable}: Store \texttt{needles.length} in a variable to eliminate repeated evaluation of the same expression. This also improves memory efficiency.
    
    \item \textbf{Optimize \texttt{haystack.split()}}: Move the following statement outside of the loop: \texttt{String[] words = haystack.split("[ \textbackslash"\textbackslash'\textbackslash t\textbackslash n\textbackslash b\textbackslash f\textbackslash r]", 0);}. The \texttt{words} array does not change with each iteration, so splitting the \texttt{haystack} only once is more efficient.
    
    \item \textbf{Improve readability}: In the third for loop, use \texttt{k} as the iterator variable, as \texttt{i} and \texttt{j} are already in use in nearby loops.
    
    \item \textbf{Consider using a \texttt{HashMap}}: To store the frequency of each needle, consider using a \texttt{HashMap}. This provides a more flexible and efficient way to handle the counting, especially if you want to return key-value pairs or add further functionality later.
\end{itemize}

\subsection*{Suggested code}

Here is the revised code, which incorporates the suggested improvements for efficiency and clarity:

\begin{lstlisting}[style=java]
public class NeedlesInHaystack {
    public static void findNeedles(String haystack, String[] needles) {
        // Store the length of the needles array in a variable
        int needlesLength = needles.length;
        
        // Split the haystack string once
        String[] words = haystack.split("[ \"\'\t\n\b\f\r]", 0);
        
        // Create an array to store the frequency counts
        int[] countArray = new int[needlesLength];

        // Iterate through the needles and count their occurrences
        for (int i = 0; i < needlesLength; i++) {
            for (int j = 0; j < words.length; j++) {
                if (words[j].compareTo(needles[i]) == 0) {
                    countArray[i]++;
                }
            }
        }

        // Print the results
        for (int k = 0; k < needlesLength; k++) {
            System.out.println(needles[k] + ": " + countArray[k]);
        }
    }

    public static void main(String[] args) {
        // Hard-coded values for demonstration
        // Ideally, values should be received from standard input
        String haystack = "Google Cloud provides APIs to use Google's ML/AI capabilities.";
        String[] needles = {"Google", "API", "documentation", "AWS", "ML/AI"};
        findNeedles(haystack, needles);
    }
}
\end{lstlisting}

\subsection*{Sample output}

\begin{lstlisting}
Google: 2
API: 0
documentation: 0
AWS: 0
ML/AI: 1
\end{lstlisting}

\end{document}
